\chapter{How to Answer a Data Science Case Study}
\companion{InfoEdge DS interview experience: Round 4 case studies (at 9:07)}{https://youtu.be/nWyBQhGnHaM?t=547}

Round 4 (45--60 minutes, a senior data scientist): ``A case study based on a real world data science problem to assess the candidate's talent in
applying machine learning theory to practice. The candidate shall be assessed on her/his ability to understand data related challenges, modelling
and design choices as well as model assessment skills. There is generally no single best solution. The interviewer will rather assess the
candidate on her/his thought process and on how she/he may fare in a team-working environment.'' The 2026--27 talk added: ``an AI case study:
asked how you would approach the problem, what type of data you will need'', and that they want people ``who actually want to understand
business''.

This chapter gives you a repeatable structure. Chapter 29 applies it to InfoEdge's own products.

\section{The ten-step skeleton}
Memorise it; say the steps aloud as you go, so the interviewer can follow and steer.
\begin{enumerate}
\item \textbf{Clarify the goal and the decision.} What business outcome? Who acts on the prediction, and how? (``Rank leads for 50 telecallers who make 200
  calls each per day.'') If clarifying questions are not allowed, \emph{state your assumptions} explicitly.
\item \textbf{Define success.} One primary business metric (conversions per call, applies per job, revenue), guardrails (complaints, unsubscribes, latency), and the
  offline proxy metric you will optimise (precision@k, NDCG, AUC, MAPE).
\item \textbf{Frame the ML problem.} Classification, regression, ranking, recommendation, clustering, forecasting, anomaly detection? What is one row (user,
  user--day, user--listing pair)? What is the label and over what time window? When is the prediction made?
\item \textbf{Data.} What data exist (logs, profiles, listings, CRM, call outcomes)? What is missing? How are labels obtained (natural outcomes, human labels, proxies)?
  Biases in how data were collected (only called leads have call outcomes: selection bias).
\item \textbf{Features.} Brainstorm by entity and by time window: user, item, user$\times$item, context, recency/frequency/monetary, trends, ratios; text and image
  embeddings. Check leakage for every feature: ``is this known at prediction time?''
\item \textbf{Baseline.} A rule or simple model the business already uses or would accept (most-recent-activity rule, popularity, logistic regression).
\item \textbf{Model.} Justify the choice by data type, size, latency, interpretability (GBM for tabular, two-tower for retrieval, Transformer for text). Handle
  imbalance and cold start.
\item \textbf{Offline evaluation.} Time-based split, the right metric at the operating point, segment analysis (new vs returning users, cities), calibration.
\item \textbf{Online evaluation and deployment.} A/B test with a clear hypothesis and sample size; batch vs real-time scoring; monitoring (data drift, prediction
  drift, metric decay); retraining cadence; feedback loops.
\item \textbf{Risks and iteration.} Fairness, privacy, gaming, cold start, exploration vs exploitation; next steps.
\end{enumerate}

\begin{intuition}[: why structure wins]
The interviewer is imagining you in their team next month. A structured answer shows you will not jump into modelling before knowing what decision the model
serves, and that you will catch leakage and evaluation mistakes on your own.
\end{intuition}

\section{Feature engineering: a generator of ideas}
When the interviewer ``keeps pushing until you run out of answers'' (as reported for the 99acres question), use this grid to generate features systematically.
\begin{center}\small
\begin{tabularx}{\linewidth}{l X}
\toprule
Dimension & Examples (99acres buyer intent)\\
\midrule
Recency & days since last session, since last enquiry\\
Frequency & sessions in last 7/30 days, listings viewed, searches\\
Depth / intensity & time on listing pages, photos viewed, floor plans opened, EMI calculator used, contact number revealed\\
Consistency (concentration) & share of views within the dominant price band (the price-band concentration feature that impressed the interviewer), share in top
locality, entropy of property types viewed (low entropy = focused buyer)\\
Trend & views this week vs last week; narrowing of price range over time\\
Context & device, time of day, city tier, festival season, interest rates\\
Profile & stated budget, loan pre-approval, family size, owner vs tenant\\
User $\times$ item & distance of viewed listings from workplace; match between budget and viewed prices\\
Social / market & inventory in the target locality, price trend of that locality\\
Text & search queries (``ready to move'', ``near metro''), enquiry messages\\
\bottomrule
\end{tabularx}
\end{center}
For each feature, state the hypothesis (``a narrow price band means a decided buyer'') and the leakage check.

\section{Metrics: from business to offline}
\begin{center}\small
\begin{tabularx}{\linewidth}{X X X}
\toprule
Business decision & Offline metric & Online metric\\
\midrule
Who to call with limited capacity & precision@k, expected conversions in top $k$ & conversions per call-hour\\
Rank search results & NDCG@10, MRR, recall@k for retrieval & CTR, contacts/applies per search, time to first contact\\
Recommend jobs & recall@k, NDCG, coverage, novelty & applies per user, recruiter responses\\
Flag fraud/spam & recall at fixed precision, PR AUC & fraud caught, false-positive complaints\\
Estimate salary & MAPE/RMSLE, interval coverage & user trust, engagement with salary card\\
Forecast demand & MAPE, bias & staffing efficiency\\
\bottomrule
\end{tabularx}
\end{center}

\section{A/B testing in one page}
\begin{itemize}
\item \textbf{Hypothesis}: the new lead scorer increases conversions per call by at least 5\%.
\item \textbf{Randomisation unit}: caller or lead (avoid interference: two variants sharing one caller's queue contaminate each other).
\item \textbf{Sample size}: depends on baseline rate, minimum detectable effect, significance $\alpha$ (0.05) and power (0.8). For a proportion $p$ and absolute effect $\delta$,
  roughly $n\approx16p(1-p)/\delta^2$ per arm. Baseline 5\%, detect +0.5 points: $16\cdot0.0475/0.000025\approx30{,}400$ per arm.
\item \textbf{Guardrails}, \textbf{novelty effects} (run long enough), \textbf{multiple testing} (pre-register the primary metric), \textbf{peeking} (do not stop at the first
  significant day).
\end{itemize}

\section{Product and business questions}
InfoEdge is a marketplace company: Naukri (jobs), 99acres (real estate), Jeevansathi (matrimony), Shiksha (education), plus investments. Its revenue is mostly from the
paying side of each marketplace (recruiters buying job postings and Resdex database access; property developers and agents buying listings and leads). Marketplaces have
\textbf{network effects}: more recruiters attract more seekers and vice versa. Good product answers grow both sides, improve match quality (which raises willingness to pay),
and use data products (Talent Pulse) to open new revenue. A useful skeleton for ``grow revenue'': \emph{more customers} (acquisition, new segments), \emph{more revenue per
customer} (upsell, pricing, premium tiers), \emph{less churn} (better outcomes, success management), \emph{new products} (data/insight products, AI tools).

\section{Communication and team fit}
Think aloud; summarise every few minutes; ask ``shall I go deeper here or move on?''; accept pushback gracefully (``good point; that would bias the labels; I would
fix it by\dots''); do not claim certainty you do not have; quantify where possible; end with a crisp summary and next steps.

\stuck{Exponent: How to answer ML system design interview questions (search)}{\yts{machine+learning+system+design+interview+framework+example}}
\section{Framework questions}

\begin{qa}{\pyq{Round 4 format: no clarifying questions allowed} How do you structure a case answer when the interviewer says you may not ask clarifying questions?}
\oneline{State your assumptions explicitly and briefly, choose the most plausible interpretation, and keep going with the usual structure; mention how the answer would change
under a different assumption.}
\why The reported Naukri revenue question disallowed questions to test company knowledge and decisiveness. ``I will assume most revenue comes from recruiters buying
postings and database access, and that growth in seekers drives recruiter value. If instead the focus were seeker subscriptions, I would\dots''
\end{qa}

\begin{qa}{What are the most common mistakes in ML case interviews?}
\oneline{Jumping to a model before defining the decision and metric, ignoring how labels are created, missing leakage, using random splits on temporal data, choosing the wrong metric
(accuracy on imbalanced data), forgetting deployment and monitoring, and never connecting back to the business outcome.}
\end{qa}

\begin{qa}{\deep Your model will change who gets shown which jobs, which changes future training data. Why is that a problem and what do you do?}
\oneline{It creates a feedback loop: the model only observes outcomes for items it chose to show, so its training data become biased toward its own past choices and it can lock in
errors or popularity bias; counter with exploration (randomised or bandit traffic), propensity-weighted training and logging of what was shown.}
\why Position bias too: users click top results regardless. Use inverse propensity scoring or position-aware models; reserve a small random slice of traffic for unbiased evaluation.
\end{qa}

\begin{qa}{How do you choose the label for ``will this user buy a property''?}
\oneline{Define an observable event and a window tied to the decision: for example ``a site visit booked or a transaction recorded within 30 days of the prediction date''; check that
it is measured reliably, that it is available for all users (not only those called), and that the window matches the business cycle.}
\why A purchase may happen offline and never be logged; proxies (site visit, developer-confirmed lead) may be better measured. Different windows give different models; choose
by how long telecalling can usefully wait.
\end{qa}

\begin{qa}{\deep Offline evaluation looks great but you can only A/B test with 5\% of traffic. How do you de-risk?}
\oneline{Use a careful time-based offline backtest, shadow deployment (score live traffic without acting), a small ramped A/B test with guardrails and pre-computed power, interleaving
for ranking (very sensitive with little traffic), and clear rollback criteria.}
\end{qa}

\begin{qa}{How would you explain a model's output to a non-technical stakeholder, such as a sales head?}
\oneline{Translate into the decision and its impact (``calling the top 200 leads converts 3 times more than the current list''), show a few concrete examples with the top reasons per lead,
and state limitations plainly.}
\end{qa}

\begin{qa}{Estimate the sample size for an A/B test with baseline conversion 5\% and a minimum detectable absolute lift of 0.5 percentage points.}
\oneline{About $16p(1-p)/\delta^2 = 16\times0.0475/0.005^2\approx30{,}400$ users per arm (for 5\% significance and 80\% power).}
\end{qa}

\begin{qa}{\deep When would you \emph{not} use machine learning for a business problem?}
\oneline{When a simple rule achieves nearly the same outcome, when there is not enough reliable labelled data, when the cost of errors or lack of explainability is unacceptable, or when the
decision is rare and better made by people; ML must beat a baseline by enough to justify maintenance cost.}
\end{qa}
